% Preámbulo común para las soluciones de los prácticos de Programación 3
\documentclass[11pt,a4paper]{article}

\usepackage[utf8]{inputenc}
\usepackage[T1]{fontenc}
\usepackage{lmodern}
\usepackage[spanish,es-noshorthands]{babel}
\usepackage[margin=2.3cm]{geometry}
\usepackage{amsmath,amssymb,amsthm}
\usepackage{enumitem}
\usepackage{xcolor}
\usepackage{tikz}
\usetikzlibrary{arrows.meta,positioning,shapes,calc}
\usepackage[noend]{algpseudocode}
\usepackage{algorithm}
\usepackage{booktabs}
\usepackage{tcolorbox}
\usepackage[hidelinks]{hyperref}

\floatname{algorithm}{Algoritmo}
\renewcommand{\algorithmicrequire}{\textbf{Entrada:}}
\renewcommand{\algorithmicensure}{\textbf{Salida:}}
\algrenewcommand\algorithmicwhile{\textbf{mientras}}
\algrenewcommand\algorithmicdo{\textbf{hacer}}
\algrenewcommand\algorithmicfor{\textbf{para}}
\algrenewcommand\algorithmicforall{\textbf{para todo}}
\algrenewcommand\algorithmicif{\textbf{si}}
\algrenewcommand\algorithmicthen{\textbf{entonces}}
\algrenewcommand\algorithmicelse{\textbf{sino}}
\algrenewcommand\algorithmicreturn{\textbf{retornar}}
\algrenewcommand\algorithmicfunction{\textbf{función}}
\algrenewcommand\algorithmicprocedure{\textbf{procedimiento}}

\theoremstyle{plain}
\newtheorem{lema}{Lema}[section]
\newtheorem{teorema}[lema]{Teorema}
\newtheorem{proposicion}[lema]{Proposición}
\theoremstyle{definition}
\newtheorem{definicion}[lema]{Definición}
\theoremstyle{remark}
\newtheorem*{observacion}{Observación}

% Caja para el enunciado de cada ejercicio
\newtcolorbox{enunciado}{colback=gray!6,colframe=gray!50,boxrule=0.4pt,arc=2pt,left=6pt,right=6pt,top=4pt,bottom=4pt,fontupper=\small}
% Nota que vincula el ejercicio con parciales/exámenes
\newcommand{\examref}[1]{\par\noindent{\small\color{blue!50!black}\textbf{Relación con parciales/exámenes:} #1}\par\medskip}

\newcommand{\NP}{\mathsf{NP}}
\newcommand{\Pclass}{\mathsf{P}}
\newcommand{\Oh}{\mathcal{O}}

\setlength{\parskip}{0.4em}
\setlength{\parindent}{0pt}


\title{Práctico 1 --- Emparejamiento estable}
\author{Programación 3 -- FING, UdelaR}
\date{2026}

\begin{document}
\maketitle

\paragraph{Notación y resultados usados.} Trabajamos con dos conjuntos $M$ y $W$ de $n$
elementos cada uno, cada persona con una lista de preferencias estricta y completa sobre el
otro conjunto. Un \emph{emparejamiento perfecto} $S$ es \emph{estable} si no tiene
inestabilidades, es decir, no existen $(m,w)\in S$ y $(m',w')\in S$ tales que $m$ prefiere
$w'$ a $w$ y $w'$ prefiere $m$ a $m'$. Decimos que $w$ es una \emph{pareja válida} de $m$
(y viceversa) si existe algún emparejamiento estable que contiene a $(m,w)$. Usamos los
siguientes resultados del libro (Kleinberg \& Tardos, Cap.~1):
\begin{itemize}
  \item (1.3) El algoritmo de Gale--Shapley termina luego de a lo sumo $n^2$ iteraciones.
  \item (1.4)--(1.6) Durante la ejecución, una $w$ que recibió una propuesta queda emparejada
        para siempre y sus parejas sucesivas son cada vez mejores (según $w$); las propuestas de
        cada $m$ siguen su orden de preferencia. El resultado es un emparejamiento perfecto y estable.
  \item (1.7) Toda ejecución de Gale--Shapley con $M$ como proponente devuelve el mismo
        emparejamiento $S^*$, en el que cada $m\in M$ queda con su \emph{mejor} pareja válida.
  \item (1.8) En $S^*$ cada $w\in W$ queda con su \emph{peor} pareja válida.
\end{itemize}

%=====================================================================
\section{Ejercicio 1}

\begin{enunciado}
\textbf{(a)} (K\&T, Ex.~1.2) Indique si la siguiente afirmación es verdadera o falsa. Si es
verdadera, dé una breve explicación; si es falsa, dé un contraejemplo.
\emph{Consideramos una instancia del problema de emparejamiento estable en la que existe un
par $(m,w)$ tal que $m$ está primero en la lista de preferencias de $w$ y $w$ está primero en
la lista de preferencias de $m$. Entonces el par $(m,w)$ pertenece a toda solución $S$ para
esta instancia del problema.}

\textbf{(b)} Dé una instancia (mínima) del problema de emparejamiento estable para la cual hay
más de un emparejamiento estable. Justifique cómo se obtiene.
\end{enunciado}

\examref{En el Examen de diciembre 2017, Pregunta~3, se pregunta por una afirmación
parecida pero distinta (\emph{en toda instancia hay un emparejamiento estable que contiene un
par que se rankea mutuamente primero}), que es falsa: el contraejemplo oficial es la misma
instancia que usamos en la parte (b).}

\subsection*{Parte (a)}
La afirmación es \textbf{verdadera}.

Sea $S$ un emparejamiento estable cualquiera y supongamos, por absurdo, que $(m,w)\notin S$.
Como $S$ es perfecto, $m$ está emparejado en $S$ con alguna $w'\neq w$ y $w$ está emparejada
en $S$ con algún $m'\neq m$; es decir, $(m,w')\in S$ y $(m',w)\in S$. Entonces:
\begin{itemize}
  \item $m$ prefiere $w$ a $w'$, porque $w$ es la primera de su lista y $w'\neq w$;
  \item $w$ prefiere $m$ a $m'$, porque $m$ es el primero de su lista y $m'\neq m$.
\end{itemize}
Por lo tanto el par $(m,w)$ es una inestabilidad de $S$ con respecto a $(m,w')$ y
$(m',w)$, lo que contradice que $S$ sea estable. Concluimos que $(m,w)\in S$ para toda
solución $S$.

\subsection*{Parte (b)}
Con $n=1$ hay un único emparejamiento perfecto, así que una instancia con más de un
emparejamiento estable debe tener $n\ge 2$. Tomamos $n=2$, $M=\{m_1,m_2\}$,
$W=\{w_1,w_2\}$ y las preferencias
\begin{center}
\begin{tabular}{l@{\qquad}l}
$m_1$: $w_1 \succ w_2$ & $w_1$: $m_2 \succ m_1$\\
$m_2$: $w_2 \succ w_1$ & $w_2$: $m_1 \succ m_2$
\end{tabular}
\end{center}
Hay solo dos emparejamientos perfectos y \textbf{ambos son estables}:
\begin{itemize}
  \item $S_1=\{(m_1,w_1),(m_2,w_2)\}$. Se obtiene ejecutando Gale--Shapley con $M$ como
        proponente: $m_1$ se propone a $w_1$ (libre, se forma $(m_1,w_1)$) y $m_2$ se propone a
        $w_2$ (libre, se forma $(m_2,w_2)$); no quedan hombres libres. Es estable por la
        corrección de Gale--Shapley; directamente: cada $m$ tiene a su primera opción, así que
        ningún $m$ prefiere a otra mujer y no puede haber inestabilidad (una inestabilidad requiere
        un $m$ que prefiera otra $w'$ a su pareja).
  \item $S_2=\{(m_1,w_2),(m_2,w_1)\}$. Se obtiene ejecutando Gale--Shapley con $W$ como
        proponente: $w_1$ se propone a $m_2$ y $w_2$ a $m_1$, ambos libres. Es estable por el
        argumento simétrico: cada mujer tiene a su primera opción.
\end{itemize}
Por lo tanto esta instancia, de tamaño mínimo $n=2$, tiene dos emparejamientos estables. En
general, una forma sistemática de detectar más de un emparejamiento estable es ejecutar
Gale--Shapley con cada conjunto como proponente: si los resultados difieren, por (1.6) son dos
emparejamientos estables distintos.

%=====================================================================
\section{Ejercicio 2}

\begin{enunciado}
(Primer parcial de 2019.) Considere la instancia con $M=\{m_1,m_2,m_3,m_4\}$,
$W=\{w_1,w_2,w_3,w_4\}$ y preferencias (de mayor a menor):
\begin{center}
\begin{tabular}{l@{\quad}llll@{\qquad\qquad}l@{\quad}llll}
$m_1$ & $w_1$ & $w_2$ & $w_3$ & $w_4$ & $w_1$ & $m_4$ & $m_3$ & $m_1$ & $m_2$\\
$m_2$ & $w_1$ & $w_4$ & $w_3$ & $w_2$ & $w_2$ & $m_2$ & $m_4$ & $m_1$ & $m_3$\\
$m_3$ & $w_2$ & $w_1$ & $w_3$ & $w_4$ & $w_3$ & $m_4$ & $m_1$ & $m_2$ & $m_3$\\
$m_4$ & $w_4$ & $w_2$ & $w_3$ & $w_1$ & $w_4$ & $m_3$ & $m_2$ & $m_1$ & $m_4$
\end{tabular}
\end{center}
\textbf{(a)} Presente una sucesión de formaciones y separaciones de parejas que se realizan a
lo largo de una ejecución del algoritmo de Gale--Shapley.
\textbf{(b)} Repita la parte (a) invirtiendo los roles de $M$ y $W$.
\textbf{(c)} El emparejamiento obtenido en ambas partes es el mismo. Muestre que, en esta
instancia, para toda persona su mejor pareja válida coincide con su peor pareja válida. ¿Qué
resultados del libro está aplicando?
\textbf{(d)} Concluya que para esta instancia existe una \textbf{única} solución.
\end{enunciado}

\examref{Es exactamente el Ejercicio~1 del Primer parcial 2019 (la parte (b) del parcial
agrupa las partes (b)--(d) de este ejercicio). Relacionado: Primer parcial 2025, Ej.~1(f)
(decidir si una instancia tiene un único emparejamiento estable comparando las salidas de
Gale--Shapley con cada conjunto como proponente).}

\subsection*{Parte (a): $M$ propone}
Elegimos en cada iteración un hombre libre (cualquiera sirve, por (1.7) el resultado no
depende de la elección). La tabla registra cada propuesta y su efecto.
\begin{center}
\begin{tabular}{cllp{7.5cm}}
\toprule
Paso & Propone & a & Resultado \\
\midrule
1 & $m_1$ & $w_1$ & $w_1$ libre: \textbf{se forma} $(m_1,w_1)$.\\
2 & $m_2$ & $w_1$ & $w_1$ prefiere $m_1$ (3.º) a $m_2$ (4.º): rechaza a $m_2$.\\
3 & $m_3$ & $w_2$ & $w_2$ libre: \textbf{se forma} $(m_3,w_2)$.\\
4 & $m_4$ & $w_4$ & $w_4$ libre: \textbf{se forma} $(m_4,w_4)$.\\
5 & $m_2$ & $w_4$ & $w_4$ prefiere $m_2$ (2.º) a $m_4$ (4.º): \textbf{se separa} $(m_4,w_4)$ y \textbf{se forma} $(m_2,w_4)$.\\
6 & $m_4$ & $w_2$ & $w_2$ prefiere $m_4$ (2.º) a $m_3$ (4.º): \textbf{se separa} $(m_3,w_2)$ y \textbf{se forma} $(m_4,w_2)$.\\
7 & $m_3$ & $w_1$ & $w_1$ prefiere $m_3$ (2.º) a $m_1$ (3.º): \textbf{se separa} $(m_1,w_1)$ y \textbf{se forma} $(m_3,w_1)$.\\
8 & $m_1$ & $w_2$ & $w_2$ prefiere $m_4$ (2.º) a $m_1$ (3.º): rechaza a $m_1$.\\
9 & $m_1$ & $w_3$ & $w_3$ libre: \textbf{se forma} $(m_1,w_3)$.\\
\bottomrule
\end{tabular}
\end{center}
No quedan hombres libres y el algoritmo devuelve
\[
S_M=\{(m_1,w_3),(m_2,w_4),(m_3,w_1),(m_4,w_2)\}.
\]
En resumen: se forman $(m_1,w_1)$, $(m_3,w_2)$, $(m_4,w_4)$; se separa $(m_4,w_4)$ y se forma
$(m_2,w_4)$; se separa $(m_3,w_2)$ y se forma $(m_4,w_2)$; se separa $(m_1,w_1)$ y se forma
$(m_3,w_1)$; se forma $(m_1,w_3)$.

\subsection*{Parte (b): $W$ propone}
\begin{center}
\begin{tabular}{cllp{7.5cm}}
\toprule
Paso & Propone & a & Resultado \\
\midrule
1 & $w_1$ & $m_4$ & $m_4$ libre: \textbf{se forma} $(w_1,m_4)$.\\
2 & $w_2$ & $m_2$ & $m_2$ libre: \textbf{se forma} $(w_2,m_2)$.\\
3 & $w_3$ & $m_4$ & $m_4$ prefiere $w_3$ (3.ª) a $w_1$ (4.ª): \textbf{se separa} $(w_1,m_4)$ y \textbf{se forma} $(w_3,m_4)$.\\
4 & $w_4$ & $m_3$ & $m_3$ libre: \textbf{se forma} $(w_4,m_3)$.\\
5 & $w_1$ & $m_3$ & $m_3$ prefiere $w_1$ (2.ª) a $w_4$ (4.ª): \textbf{se separa} $(w_4,m_3)$ y \textbf{se forma} $(w_1,m_3)$.\\
6 & $w_4$ & $m_2$ & $m_2$ prefiere $w_4$ (2.ª) a $w_2$ (4.ª): \textbf{se separa} $(w_2,m_2)$ y \textbf{se forma} $(w_4,m_2)$.\\
7 & $w_2$ & $m_4$ & $m_4$ prefiere $w_2$ (2.ª) a $w_3$ (3.ª): \textbf{se separa} $(w_3,m_4)$ y \textbf{se forma} $(w_2,m_4)$.\\
8 & $w_3$ & $m_1$ & $m_1$ libre: \textbf{se forma} $(w_3,m_1)$.\\
\bottomrule
\end{tabular}
\end{center}
No quedan mujeres libres y el resultado es
\[
S_W=\{(w_1,m_3),(w_2,m_4),(w_3,m_1),(w_4,m_2)\},
\]
que, salvo el orden de los componentes de cada par, es el mismo emparejamiento que $S_M$.

\subsection*{Parte (c)}
Aplicamos los resultados (1.7) y (1.8) del libro, en las dos versiones que se obtienen al
intercambiar los roles de $M$ y $W$ (los resultados son simétricos, porque el problema lo es):
\begin{itemize}
  \item Con $M$ proponente, por (1.7) en $S_M$ cada $m$ queda con su \emph{mejor} pareja
        válida, y por (1.8) cada $w$ queda con su \emph{peor} pareja válida.
  \item Con $W$ proponente, por (1.7) (con roles intercambiados) en $S_W$ cada $w$ queda con su
        \emph{mejor} pareja válida, y por (1.8) cada $m$ queda con su \emph{peor} pareja válida.
\end{itemize}
Como $S_M=S_W$, para cada $m\in M$ su mejor pareja válida (su pareja en $S_M$) es igual a su
peor pareja válida (su pareja en $S_W$); análogamente para cada $w\in W$. Concretamente:
\[
\begin{array}{c|cccc|cccc}
 & m_1 & m_2 & m_3 & m_4 & w_1 & w_2 & w_3 & w_4\\\hline
\text{mejor = peor pareja válida} & w_3 & w_4 & w_1 & w_2 & m_3 & m_4 & m_1 & m_2
\end{array}
\]

\subsection*{Parte (d)}
Sea $S$ un emparejamiento estable cualquiera y sea $(m,w)\in S$. Por definición, $w$ es pareja
válida de $m$, de modo que, en la lista de $m$, $w$ no puede estar antes que la mejor pareja
válida de $m$ ni después que su peor pareja válida. Como por la parte (c) ambas coinciden con
$S_M(m)$, se tiene $w=S_M(m)$. Esto vale para todo $m$, así que $S=S_M$. Por lo tanto la
instancia tiene una \textbf{única} solución, $S_M=S_W$.

%=====================================================================
\section{Ejercicio 3}

\begin{enunciado}
(K\&T, Ex.~1.4) Hay un conjunto $H$ de $m$ hospitales y un conjunto $E$ de $n$ estudiantes;
los hospitales tienen listas de preferencias (completas) sobre $E$ y los estudiantes sobre $H$.
Cada hospital $h$ tiene un cupo $C(h)$ y se cumple $\sum_{h\in H} C(h)\le n$. Se quiere
asignar estudiantes a hospitales, de modo que cada estudiante quede asignado a lo sumo a un
hospital y todas las plazas queden ocupadas, sin inestabilidades:
\begin{itemize}
  \item \emph{Tipo 1:} existen $e,e'$ y $h$ tales que $e$ está asignado a $h$, $e'$ está libre y
        $h$ prefiere a $e'$ antes que a $e$.
  \item \emph{Tipo 2:} existen $e,e'$ y $h,h'$ tales que $e$ está asignado a $h$, $e'$ está
        asignado a $h'$, $h$ prefiere a $e'$ antes que a $e$ y $e'$ prefiere a $h$ antes que a $h'$.
\end{itemize}
Demuestre que para toda instancia existe una asignación estable y presente un algoritmo para
obtenerla. \emph{Sugerencia:} adapte Gale--Shapley con los hospitales como proponentes y muestre
(1) que termina, (2) que al terminar todas las plazas están ocupadas, (3) que no hay
inestabilidades de tipo 1 y (4) que no hay inestabilidades de tipo 2.
\end{enunciado}

\examref{El Primer parcial 2022, Ej.~1 (asignación de jugadores a servidores con cupo
$|J|/|S|$) es una variante de este ejercicio; la solución oficial usa la misma adaptación de
Gale--Shapley con los servidores (análogos a los hospitales) como proponentes, y acota las
iteraciones por la cantidad de pares (servidor, jugador).}

\subsection*{Algoritmo}
Cada hospital lleva la cuenta de cuántos estudiantes tiene asignados, $a(h)$. Decimos que $h$
\emph{tiene plazas libres} si $a(h)<C(h)$.

\begin{algorithm}[H]
\caption{Gale--Shapley para hospitales y estudiantes}
\begin{algorithmic}[1]
\State Inicialmente todo estudiante está libre y $a(h)=0$ para todo $h\in H$
\While{existe $h\in H$ con $a(h)<C(h)$ que no se propuso a todos los estudiantes}
  \State Sea $e$ el estudiante de mayor preferencia para $h$ al que $h$ no se propuso
  \If{$e$ está libre}
    \State Asignar $e$ a $h$; \ $a(h)\gets a(h)+1$
  \Else
    \State Sea $h'$ el hospital al que está asignado $e$
    \If{$e$ prefiere $h$ antes que $h'$}
      \State Desasignar $e$ de $h'$; \ $a(h')\gets a(h')-1$
      \State Asignar $e$ a $h$; \ $a(h)\gets a(h)+1$
    \Else
      \State $e$ rechaza a $h$
    \EndIf
  \EndIf
\EndWhile
\State \Return la asignación
\end{algorithmic}
\end{algorithm}

\paragraph{Invariantes.} Las siguientes propiedades se verifican trivialmente por inducción en
las iteraciones, mirando los pasos del algoritmo:
\begin{enumerate}[label=(I\arabic*)]
  \item Cada estudiante está asignado a lo sumo a un hospital, y para todo $h$ el valor $a(h)$
        es la cantidad de estudiantes asignados a $h$, con $0\le a(h)\le C(h)$. (Un hospital
        recibe un estudiante sólo en los pasos 5 y 10, cuando se propone, y sólo se propone si
        $a(h)<C(h)$; en el paso 9 un estudiante deja $h'$ y pasa a $h$, sin quedar en dos
        hospitales.)
  \item Una vez que un estudiante recibe una propuesta, queda asignado a algún hospital hasta el
        final (en los pasos 5 y 10 queda asignado y nunca se lo deja libre: en el paso 9 se lo
        desasigna sólo para asignarlo en el paso 10).
  \item Los hospitales a los que está asignado un estudiante son cada vez mejores según sus
        preferencias (sólo cambia en el paso 10, y sólo a un hospital que prefiere).
  \item Cada hospital se propone a los estudiantes en orden decreciente de preferencia, y nunca
        dos veces al mismo (paso 3).
  \item Si al final $e$ está asignado a $h$, entonces $h$ se propuso a $e$ (la única forma de
        quedar asignado a $h$ es aceptar una propuesta de $h$).
\end{enumerate}

\subsection*{1. El algoritmo termina}
En cada iteración del ciclo un hospital $h$ se propone a un estudiante al que no se había
propuesto antes (paso 3 e (I4)). Por lo tanto el conjunto de pares $(h,e)$ tales que $h$ se
propuso a $e$ crece estrictamente en cada iteración, y como está contenido en $H\times E$, que
tiene $mn$ elementos, hay a lo sumo $mn$ iteraciones. El algoritmo termina.

Además, con las estructuras de la implementación estándar de Gale--Shapley (una lista de
hospitales con plazas libres, un índice al próximo estudiante en la lista de cada hospital, el
hospital actual de cada estudiante y una matriz \texttt{ranking}$[e][h]$ con la posición de $h$
en la lista de $e$) cada iteración cuesta $O(1)$, y el algoritmo corre en tiempo $O(mn)$, sin
contar el preprocesamiento de \texttt{ranking}, que también es $O(mn)$.

\subsection*{2. Al terminar, todas las plazas están ocupadas}
Supongamos por absurdo que al terminar existe $h$ con $a(h)<C(h)$. Como la condición del ciclo
es falsa, $h$ se propuso a todos los estudiantes. Por (I2), todos los $n$ estudiantes recibieron
una propuesta y están asignados. Por (I1) cada estudiante ocupa exactamente una plaza, así que
la cantidad de plazas ocupadas es
\[
n=\sum_{g\in H} a(g) < \sum_{g\in H} C(g)\le n,
\]
donde la desigualdad estricta se debe a que $a(h)<C(h)$ y $a(g)\le C(g)$ para todo $g$. Esto es
una contradicción. Luego todas las plazas están ocupadas, y junto con (I1) el resultado es una
asignación en la que cada estudiante está asignado a lo sumo a un hospital y todos los
hospitales llenan sus cupos.

\subsection*{3. No hay inestabilidades de tipo 1}
Supongamos por absurdo que en la asignación final $e$ está asignado a $h$, $e'$ está libre y
$h$ prefiere $e'$ antes que $e$. Por (I5), $h$ se propuso a $e$; por (I4), como $h$ prefiere
$e'$ a $e$, $h$ se propuso a $e'$ antes de proponerse a $e$. Pero entonces, por (I2), $e'$ quedó
asignado desde ese momento hasta el final, lo que contradice que $e'$ esté libre.

\subsection*{4. No hay inestabilidades de tipo 2}
Supongamos por absurdo que en la asignación final $e$ está asignado a $h$, $e'$ está asignado a
$h'$, $h$ prefiere $e'$ antes que $e$ y $e'$ prefiere $h$ antes que $h'$. Notar que $h'\ne h$,
porque $e'$ prefiere estrictamente $h$ a $h'$.

Por (I5) $h$ se propuso a $e$ y, por (I4), como prefiere $e'$, antes se propuso a $e'$. Inmediatamente después de esa propuesta, $e'$ está asignado a un hospital $h''$ tal que
$h''=h$ (si aceptó la propuesta) o $e'$ prefiere $h''$ antes que $h$ (si la rechazó, porque sólo
se rechaza cuando el hospital actual es preferido). Por (I3), el hospital final $h'$ de $e'$ es
$h''$ o uno que $e'$ prefiere a $h''$. En cualquier caso, $h'=h$ o $e'$ prefiere $h'$ antes que
$h$. Como $h'\ne h$, concluimos que $e'$ prefiere $h'$ antes que $h$, lo que contradice la
hipótesis.

\subsection*{Conclusión}
Para toda instancia el algoritmo termina (parte 1) y devuelve una asignación en la que cada
estudiante está en a lo sumo un hospital y todas las plazas están ocupadas (parte 2), sin
inestabilidades de tipo 1 (parte 3) ni de tipo 2 (parte 4); es decir, una asignación estable.
En particular, toda instancia tiene al menos una asignación estable.

\end{document}
